\section{The arithmetic window and the polynomial bound at level two}
\label{sec:arithmetic-translation}

Two statements involving arithmetic lengths occur close to the conclusion in
the source. The proof of nonvanishing for all characters uses the window
assertion denoted by \(\mathrm{VWin}\). A separate theorem in the module
\texttt{Small\allowbreak{}Class\allowbreak{}Numbers.\allowbreak{}Cor\allowbreak{}GKOf} applies a general estimate to the trivial
character modulo two. We give the closing deductions of both statements and
specify the results imported at each deduction. Merely importing a module does
not make every theorem in that module a premise of the theorem being proved.

\subsection{The four counts in the window argument}

Here is the window with its quantifiers, from \texttt{Exists\allowbreak{}Page\allowbreak{}N}. For a
number field \(E\), a prime \(\ell\), and a positive twist \(k\), let
\(s_{\ell,E}(k)\) be the \(\mathbb Z_\ell\)-length of the kernel of the
localization map in degree one for
\((\mathbb Q_\ell/\mathbb Z_\ell)(-k)\). The global group allows
ramification at the primes above \(\ell\); local restriction is taken at
each of those primes. Formally the length is defined for a choice of
algebraic closure of \(E\) and local data \(D_\ell\) at each prime.
The assertion \(\mathrm{VVle}(E,t,B)\) quantifies over every such closure
and every family \(D_\ell\). It requires finiteness for every prime and
every \(1\le k<t\), and, for every finite set \(S\) of primes, the inequality
\[
  \sum_{\ell\in S}\log\ell\sum_{1\le k<t}s_{\ell,E}(k)\ \le B.
\]
Shrinking \(t\) and increasing \(B\) preserves the assertion: the
weights and omitted lengths are nonnegative, and finiteness is inherited.

Fix a primitive character \(\chi\) of level \(m\), a real number
\(0<\varepsilon<1/52\), and any function \(a:\mathbb R\to\mathbb N\)
such that
\[
   x^{2+8\varepsilon}\le 2^{a(x)}\le 2x^{2+8\varepsilon}
       \qquad(x\ge 1).
\]
There are \(J\ge1\) and \(X_0\) such that the following statement holds
for every \(X\ge X_0\). Write \(U=X^{1/J}\) and take any primes
\(q\equiv3\bmod 8,\ d\equiv7\bmod16\), with \(q\mid d+1\), such that
\[
 \begin{split}
  \tfrac12U^{\varepsilon/5}&<\log q\le U^{\varepsilon/5},\\
  \tfrac12U^{\varepsilon/5}(\log U)^2&<\log d
         \le U^{\varepsilon/5}(\log U)^2,
 \end{split}
\]
and such that \(p\mid d+1\) for every odd prime \(p\le U^{\varepsilon/5}\).
There exists \(Y\ge U\) with \(Y^2\le X\) such that, for every
\(Y\le x\le Y^2\) and every realization of the following field \(E\),
\[
   \mathrm{VVle}\!\left(E,\lfloor x^{401}\rfloor,\,
                      \sqrt{\exp(x^\varepsilon)}\right).
\]
The field condition means that \(E\) embeds in a cyclotomic field at a
level divisible by \(2^{a(x)+2},m,q,d\), and, for a fixed primitive
root at that level, the automorphism given by exponent \(b\) fixes
the image of \(E\) pointwise precisely when
\[
 b\equiv\pm1\bmod 2^{a(x)+2},\qquad
 \chi(b)=1,\qquad
 \left(\frac bq\right)=\left(\frac bd\right)=1.
\]
In particular \(q,d\) are fixed throughout the interval for \(x\), while
the two-power exponent \(a(x)\) can vary there.

Let \(H_N\) denote the collection of local representation-theoretic,
automorphic, Galois-cohomological and complex-multiplication results called
\texttt{Literature\allowbreak{}Rank\allowbreak{}NHypothesis}. In the proof of
\texttt{vwin\_\allowbreak{}of\_\allowbreak{}row}, the input \(h\) is the Galois compatibility assertion
for the specified local parameter relation. The open-image record is formed
from \(h\). The local inputs are obtained from this record, the automorphic
inputs from this record together with base change and \(h\), and the Galois
cohomology and complex-multiplication inputs from the same record. These four
constructions give \(H_N\). No auxiliary field has yet been specialized to
discriminant \(-3\).

Apply \texttt{published\_\allowbreak{}hypotheses\_\allowbreak{}from\_\allowbreak{}patching\_\allowbreak{}alt} to \(H_N\) and the
weak descent theorem \texttt{desc\allowbreak{}Weak}, and call the resulting data \(H'\).
The size lemma called \(K\) supplies a real constant \(C_K\) from the
rank-\(N\) part of \(H'\). The endpoint theorem
\texttt{ends0\_\allowbreak{}alt} supplies two assertions. The first, together with
the rank-\(N\) inputs, produces unramified-coefficient automorphy. With
the two-dimensional Kato input, \(H'\), and the size lemma \(K\), this
produces the count in case I. The two counts in cases II and III are
obtained together by applying
\texttt{count\allowbreak{}II\_\allowbreak{}and\_\allowbreak{}count\allowbreak{}III\_\allowbreak{}of\_\allowbreak{}prints\allowbreak{}H\_\allowbreak{}of\_\allowbreak{}sp4\allowbreak{}In}
to \(H'\), \(K\), and its four-dimensional symplectic input. The second
endpoint, with \(H'\) and \(K\), gives the count in case V through the
reduced unitary lift. These are four distinct imported count theorems.

The last step uses \texttt{of\_\allowbreak{}four\_\allowbreak{}alt}. Its hypotheses consist of
\(H_N\), the existence of \(C_K\), and counts I, II, III, V. Apply
\texttt{thm\allowbreak{}Bab\_\allowbreak{}of\_\allowbreak{}lemma\allowbreak{}K} to the global duality input, the two local
cohomology inputs, the Ferrero--Washington input, and the size lemma.
Specialize the resulting auxiliary estimate to \(401\). The minus-eigenspace bound is supplied by \texttt{thm\allowbreak{}Bc\allowbreak{}Free\_\allowbreak{}of} applied to \texttt{ML.\allowbreak{}iv}. On the positive side the exceptional-prime sum uses \texttt{count\allowbreak{}IV} and \texttt{count\allowbreak{}Vs} alongside II and V. Case Z is supplied by \texttt{count\allowbreak{}Z} applied to \texttt{prop\allowbreak{}Z}. The theorem \texttt{vwin\_\allowbreak{}of\_\allowbreak{}counts} now applies to these
inputs and the four counts, yielding \(\mathrm{VWin}\). This describes the
use of the four estimates; their arithmetic proofs are dependencies of this
deduction. The source loci are \texttt{Vwin\allowbreak{}Of\allowbreak{}Row},
\texttt{Top\allowbreak{}Main\allowbreak{}Zj}, \texttt{Window\allowbreak{}Bound\allowbreak{}Of\allowbreak{}Literature\allowbreak{}Level}, and
\texttt{HVCore}, all in the \texttt{Dense\allowbreak{}Threshold\allowbreak{}Cases} directory.

\subsection{Character sums in the four counts}

The sets of characters and the range of summation are part of the
estimates. Fix the target \(\chi\). A host \(h\) determines its character
modulus \(N\), the two primes \(q_1,d\), and the subgroup \(H_h\) of
\((\mathbb Z/N\mathbb Z)^\times\) fixing its abelian field.
For coefficient ring \(\mathcal O_\ell\) and twist \(k\), let
\(\mathcal X_{\ell,k}\) be the characters modulo \(N\), valued in
\(\mathcal O_\ell\), trivial on \(H_h\), and with
\(\psi(-1)=(-1)^k\). Let \(b_{\ell,k}(\psi)\) be the strict length for
the character \(\psi^{-1}\chi_\ell^{-k}\), with strict local conditions
at the primes dividing \(\ell N\); put
\(v_{\ell,k}(\psi)=b_{\ell,k}(\psi)/r_\ell\) when the length is finite.
Write \(B_\chi\) for the exceptional primes of the target. Define
\(\psi^{(0)}_\ell\) by removing the \(\ell\)-power-order part of the
character and then removing its component at the prime \(\ell\) in
the level. The cases which occur in the count statements are
\[
\begin{array}{cl}
\mathrm I&
 \ell\notin B_\chi,\quad \ell-1\nmid 2k;\\
\mathrm Z&
 \ell\notin B_\chi,\quad\ell-1\mid 2k,\quad\psi=1;\\
\mathrm {II}&
 k\ge2,\quad 2\mid\operatorname{ord}(\psi_\ell^{(0)}),\quad
    (\ell\in B_\chi\ \text{or}\ \ell-1\mid 2k);\\
\mathrm {III}&
 \ell\notin B_\chi,\quad\ell-1\mid 2k,\quad
   \operatorname{ord}(\psi)\ \text{odd},\quad\psi\ne1;\\
\mathrm {IV}&
 \ell\in B_\chi,\quad k\ge2,\quad
   \operatorname{ord}(\psi_\ell^{(0)})\ \text{odd};\\
\mathrm {V}&\ell\in B_\chi,\quad k=1.
\end{array}
\]
The host-size condition at cutoff \(t\) requires that both \(q_1,d\)
exceed \(2t+1\) and all five host parameters \(q,o,p_K,h_*,p_1\)
in \texttt{Host.\allowbreak{}Big}. In case I there exist an integer \(C\) and test
primes \(r_1,r_2\), depending on the target, such that whenever the
host meets this condition and both auxiliary primes exceed both test
primes, for each \(1\le k<t\), every choice of coefficient rings,
and every finite set \(S\) of odd coefficient primes in case I,
all the relevant lengths are finite and
\[
     \sum_{\ell\in S}\log\ell
             \sum_{\psi\in\mathcal X_{\ell,k}}v_{\ell,k}(\psi)
            \le (Nt)^C .
\]
In II and III the assertion is pointwise. For each target there is an
integer \(C\) such that at all these hosts, all \(1\le k<t\),
all odd primes, all coefficient choices and all
\(\psi\in\mathcal X_{\ell,k}\) in the specified case,
\[
      b_{\ell,k}(\psi)<\infty,\qquad
      (\log\ell)\,v_{\ell,k}(\psi)\le(\ell Nk)^C .
\]
No additional condition involving \(r_1,r_2\) occurs in the pointwise
statement. Case V has this shape on the branch
\(R_1(\ell,\psi)\le 2^{c_2(\psi)}\) of the host cutoff; the source
assigns the other branch a separate name \texttt{Count\allowbreak{}Vs}. Case Z
uses the vanishing for the trivial character. Case IV uses a
cyclotomic-layer estimate summed over twists. In particular the
pointwise estimate does not authorize summing an infinite set of primes.
The finite-prime-set quantifier in \(\mathrm{VVle}\) is supplied after
the separate estimates have been assembled.

\subsection{The exceptional primes and the sum over the other primes}

The exceptional primes can be computed before any window is chosen.
Write \(o=\operatorname{ord}(\chi)\). Let \(p_K\) be the least prime
\(p\ge7\) congruent to \(3\bmod4\) and prime to \(m\), let
\(h_*\) be the class number of \(\mathbb Q(\sqrt{-p_K})\), and
let \(p_1\) be the least prime coprime to \(2mp_K\) and congruent
to one modulo the odd part of \(o\). The definition in \texttt{Stmt\allowbreak{}Fields}
is
\[
 B_\chi=\{3\}\ \cup\
 \{p:\ p\text{ an odd prime},\ p\mid o\,m\,p_K\,h_*\,p_1\}.
\]
For the primitive trivial target, \(o=m=1\), \(p_K=7\) and \(p_1=3\).
The field \(\mathbb Q(\sqrt{-7})\) has discriminant \(-7\).
Every ideal class has an integral representative of norm at most
\(2\sqrt7/\pi<2\) by the imaginary quadratic Minkowski bound.
Its norm is one. Thus \(h_*=1\) and \(B_{\mathbf1}=\{3,7\}\).
For a general target the product in the display can have additional
odd prime factors.

For a CM realization \(E\) of the host and local data for its
algebraic closure, write \(\lambda^+_{\ell,e}\) for the
\(\mathbb Z_\ell\)-length of the submodule of strict classes
where complex conjugation acts by \((-1)^e\), called
\texttt{plus\allowbreak{}Len}. We derive the assertion
\texttt{Sum\allowbreak{}Out} of \texttt{Height\allowbreak{}And\allowbreak{}Vanishing.\allowbreak{}Sum\allowbreak{}Out}. Take the
constants \(C_1,C_Z,C_2,C_3\) supplied by the four indicated
counts for \(\chi\), and take the two test primes supplied by count I.
For every host satisfying the size and test-prime conditions above,
every integer \(t\ge2\), and every finite set \(S\) of odd primes
disjoint from \(B_\chi\), put
\[
        M=Nt,\qquad K=\max\{C_Z,C_2,C_3\}.
\]
Here \(M\ge2\). All \(\lambda^+_{\ell,e}\) for \(\ell\in S\) and
\(1\le e<t\) are finite, and
\begin{equation}
 \sum_{\ell\in S}\log\ell\sum_{1\le e<t}\lambda^+_{\ell,e}
       \le M^{C_1+3K+4}.
 \label{eq:rh-arith-sumout}
\end{equation}
The level and \(t\) vary while the exponent and test primes remain fixed.

To prove this assertion, choose at each \(\ell\) an admissible coefficient
ring with the roots of unity required by the character comparison.
Fix \(e\) in the indicated range. The contribution from all \(\ell\in S\)
in case I to the weighted character sum is at most \(M^{C_1}\).
For the other primes, \(\ell-1\mid2e\), hence
\(\ell\le 2e+1\le2t\). There are at most \(2t\) such primes, by
injecting them into the integers from \(1\) to \(2t\). The case
decomposition in \texttt{Height\allowbreak{}And\allowbreak{}Vanishing.\allowbreak{}Cases} places each of their
characters in Z, II, or III. Since \(\ell N e\le2Mt\le M^3\), the
pointwise estimate gives
\[
    (\log\ell)v_{\ell,e}(\psi)\le M^{3K}.
\]
There are at most \(N\) characters at a given prime, by
\texttt{Card\allowbreak{}Le\allowbreak{}Level.\allowbreak{}x\allowbreak{}Card}. Consequently the entire weighted character
sum at this \(e\) is at most \(M^{C_1}+2M^{3K+1}\).
Its individual lengths are finite in Z, II and III by the pointwise
assertions, and in I by applying that count to a singleton prime.

The imported field-to-character comparison \texttt{red\allowbreak{}Prime} applies at
each prime of \(S\). Here the prime is prime to the index of the
subgroup defining the host field, by \texttt{bchi\allowbreak{}Facts}, and
\texttt{host\allowbreak{}Level} identifies the field at modulus \(N\).
With zero error term the summation lemma \texttt{Chr.\allowbreak{}sum\allowbreak{}Back} gives
finiteness on the field side and
\[
     \sum_{1\le e<t}\lambda^+_{\ell,e}
        \le \sum_{1\le e<t}\sum_{\psi\in\mathcal X_{\ell,e}}
                                      v_{\ell,e}(\psi).
\]
Multiply by the nonnegative \(\log\ell\), sum over \(S\), interchange
the finite sums, and apply the preceding estimate for each \(e\).
There are at most \(t\) twists. Since \(t\le M\), the total is at most
\[
   t(M^{C_1}+2M^{3K+1})
   \le M^{C_1+1}+2M^{3K+2}
   \le M^{C_1+3K+4}.
\]
For the last inequality, each term without its numerical coefficient
is at most \(M^{C_1+3K+2}\) and \(3\le M^2\).
This proves \eqref{eq:rh-arith-sumout}.

\subsection{Restoring the exceptional primes}

Fix a function \(j:\mathbb N\to\mathbb N\). For each
\(\ell\in B_\chi\) the source uses \texttt{Layer\allowbreak{}Le} at \(j(\ell)\):
for every number field \(L_0\) realizing the abelian-field datum
\(L_\chi(\ell)\) of the full cyclotomic field of level \(\ell m\),
and every cyclotomic layer \(L_1/L_0\) obtained by
adjoining \(\ell\)-power roots of unity and having degree
\(\ell^{j(\ell)}\), the single-prime strict-length bound
\(\mathrm{VVat}(\ell,L_1,t,B)\) holds. Explicitly, for every
algebraic closure of \(L_1\) and each choice of local data, each
\(s_{\ell,L_1}(e)\) for \(1\le e<t\) is finite and
\(\sum_{1\le e<t}s_{\ell,L_1}(e)\le B\). Assume a common \(B\ge0\)
for all the specified layers. The input \texttt{Sum\allowbreak{}In} then gives,
for each exceptional prime, a natural number
\(C_\ell=C_\ell(\chi,\ell,j(\ell))\) and
\[
       \log\ell\sum_{1\le e<t}\lambda^+_{\ell,e}
                   \le M^{C_\ell}+C_\ell B,
\]
with each length finite. The same host-size condition is used above and
below. For primes outside \(B_\chi\), write \(C_0\) for the exponent
in \eqref{eq:rh-arith-sumout}. Put
\[
       b=|B_\chi|,\quad D=\sum_{\ell\in B_\chi}C_\ell,\quad
       C=C_0+D+b+bD .
\]
We claim that for every finite set \(S\) of odd primes the left side
of \eqref{eq:rh-arith-sumout} is at most \(M^C+CB\).
Indeed, \(C_\ell\le D\). Summing the exceptional estimates and then
applying \eqref{eq:rh-arith-sumout} to \(S\setminus B_\chi\) gives
\[
     |S\cap B_\chi|(M^D+DB)+M^{C_0}.
\]
The terms independent of \(B\) total at most
\((b+1)M^{D+C_0}\), since \(M\ge2\). Now \(b+1\le2^b\le M^b\)
and \(D+C_0+b\le C\). The coefficient of \(B\) is at most \(bD\le C\).
This proves the claim. Finiteness for any given odd prime follows
by its singleton estimate on the corresponding side of the partition.
This is the proof of \texttt{Height\allowbreak{}And\allowbreak{}Vanishing.\allowbreak{}Point\allowbreak{}Arith.\allowbreak{}point}.
It makes explicit the dependence on the target and its layer indices,
and the fact that the positive-length estimate still has a layer
hypothesis at every exceptional prime.

\subsection{Recombining the two eigenspaces and the prime two}

Let \(c\) denote the action of complex conjugation on the twisted
cohomology over \(E\). Functoriality of the Galois action gives
\(c^2=1\) and its compatibility with localization preserves the strict
submodule. Write \(\lambda^-_{\ell,e}\) for the length in the
\((-1)^{e+1}\)-eigenspace. For an odd \(\ell\), the element \(2\)
is invertible in \(\mathbb Z_\ell\), because its reduction modulo
\(\ell\) is nonzero. Put \(s=(-1)^e\). The strict class \(x\) has the
decomposition
\[
        x=\frac{x+scx}{2}+\frac{x-scx}{2}.
\]
The two summands lie in the strict submodule. Applying \(c\) multiplies
the first by \(s\) and the second by \(-s\). Their eigenspaces intersect
trivially, since a member of the intersection is killed by \(2s\), a
unit. Thus addition is an isomorphism from their direct sum to the
strict submodule. Additivity of module length proves
\[
       s_{\ell,E}(e)=\lambda^+_{\ell,e}+\lambda^-_{\ell,e}.
\]
This equality also holds for extended lengths. If each signed length
is finite, their sum is finite.

Suppose that the positive and negative estimates over every finite set
of odd primes are \(B_+\) and \(B_-\), respectively, with the displayed
twist range, and that the full lengths at two are finite with
\(\sum_{1\le e<t}s_{2,E}(e)\le B_2\). For any finite set of primes
\(S\), split off \(S\cap\{2\}\). Its weighted contribution is at most
\((\log 2)B_2\), because all lengths are nonnegative. Apply the
eigenspace equality at every prime in \(S\setminus\{2\}\) and sum the
two estimates there. We obtain
\[
       \mathrm{VVle}(E,t,\,(\log 2)B_2+B_-+B_+).
\]
These are the proofs of \texttt{sha\allowbreak{}Len\_\allowbreak{}eq\_\allowbreak{}add} and
\texttt{vvle\_\allowbreak{}of\_\allowbreak{}parts} in \texttt{Length\allowbreak{}Eq\allowbreak{}Add\allowbreak{}Eigen}.

\subsection{From two congruence integers to an orbit sum}

This is the implication from \(\mathrm{PropGK}\) to \(\mathrm{CorGK}\)
proved in \texttt{Unitary\allowbreak{}Three\allowbreak{}Groups.\allowbreak{}Cor\allowbreak{}GSum\allowbreak{}Log}. Fix a number field
\(K\). The antecedent gives integers \(c_5,A\ge0\) depending on
\(K\). Given \(k\ge1,\ N\ge1\), a subgroup
\(H\le(\mathbb Z/N\mathbb Z)^\times\) of index \(u\), and a prime
\(p\nmid 30N|D_K|\), there is a nonzero integer \(n_p\) with
\[
   \log|n_p|\le c_5\varphi(u)\{N(k+1)\}^A\log p.
\]
For an algebraic closure \(\Omega\) of \(\mathbb Q\) equipped with an
embedding of \(K\), and an admissible ring \(\mathcal O_\ell\), write
\(\mathcal X_\ell(H)\) for the finite set of characters valued in this
ring whose kernel on the units modulo \(N\) is exactly \(H\).
For \(\ell\ne p\), provided at least one character in
\(\mathcal X_\ell(H)\) satisfies the good-data predicate, the
antecedent gives finiteness of all lengths in that orbit and
\[
       \sum_{\psi\in\mathcal X_\ell(H)}v_{\ell,k}(\psi)
               \le 2\operatorname{ord}_\ell(n_p).
\]
The integer \(n_p\) has already been chosen before \(\Omega,\ell\)
and its coefficient ring.

Choose distinct allowable primes \(p,p'\), a finite set \(S\) of
coefficient primes and, for each \(\ell\in S\), a witness for the
good-data predicate in its orbit. When \(S\) is empty the required sum
is zero. Otherwise such a witness gives \(k\ge1\), so we have both
\(n_p\) and \(n_{p'}\). For each \(\ell\in S\), at least one is
applicable, even when \(p\) or \(p'\) belongs to \(S\). The valuation
of the other nonzero integer is nonnegative. Thus the lengths are
finite and their sum at \(\ell\) is at most
\(2\operatorname{ord}_\ell(n_p)+2\operatorname{ord}_\ell(n_{p'})\).
For every nonzero integer \(n\), the distinct prime powers
\(\ell^{\operatorname{ord}_\ell(n)}\), \(\ell\in S\), are pairwise
coprime and each divides \(n\); their product divides \(|n|\).
Taking the logarithm gives
\(\sum_{\ell\in S}\log\ell\,\operatorname{ord}_\ell(n)\le\log|n|\).
Apply this inequality twice to obtain exactly
\[
   \sum_{\ell\in S}\log\ell
         \sum_{\psi\in\mathcal X_\ell(H)}v_{\ell,k}(\psi)
     \le 2c_5\varphi(u)\{N(k+1)\}^A(\log p+\log p').
\]
This also explains why assigning a distinct auxiliary prime to every
coefficient prime would change the content of the estimate.

\subsection{A consequence for the trivial character at level two}

Fix an algebraic closure \(\Omega\) of \(\mathbb Q\). For a rational prime
\(\ell\), let \(\mathcal O_\ell\) be a permitted coefficient ring, write
\(\chi_\ell\) for its cyclotomic character, and write
\(r_\ell\) for the ramification index of the coefficient ring.
In the notation of \texttt{Small\allowbreak{}Class\allowbreak{}Numbers.\allowbreak{}Class\allowbreak{}Number\allowbreak{}Three}, put
\[
  b_\ell(e)=
  \operatorname{length}_{\mathcal O_\ell}
  \operatorname{ShaU}\!\left(
    \operatorname{twist}_2(\ell,\chi_\ell^{-e}),\,2\ell,\,\top
  \right).
\]
The underlying discrete coefficient module is
\(\mathcal O_\ell\otimes_{\mathbb Z_\ell}(\mathbb Q_\ell/\mathbb Z_\ell)\),
acted on by the indicated character; \(\mathcal O_\ell\) is a DVR
subalgebra of \(\overline{\mathbb Q}_\ell\), finite over
\(\mathbb Z_\ell\). Here \(r_\ell=\operatorname{length}_{\mathcal O_\ell}
(\mathcal O_\ell/\ell\mathcal O_\ell)\).
The symbol \(\top\) denotes the full Galois group, and the subgroup
\(\operatorname{ShaU}\) consists of degree-one classes restricting
trivially to all inertia groups and to the decomposition groups at
places above \(2\ell\). The length is initially in
\(\mathbb N\cup\{\infty\}\). The
polynomial estimate asserts that there exist natural numbers \(C,A\),
independent of all the following choices, such that for every
\(\Omega\), every positive even integer \(e\), every choice of coefficient
rings for all primes, and every finite set \(S\) of primes satisfying
\[
   \ell\geq 5,\qquad \ell-1\nmid 2e
       \qquad(\ell\in S),
\]
each \(b_\ell(e)\) for \(\ell\in S\) is finite and
\[
     \sum_{\ell\in S}(\log\ell)\,\frac{b_\ell(e)}{r_\ell}
       \ \leq\ C(e+1)^A .
\]
The division by \(r_\ell\) occurs separately in each summand.

We give the specialization proof. Every Dirichlet character modulo two
is trivial because \((\mathbb Z/2\mathbb Z)^\times\) has one element.
The Galois character attached to this trivial character is also trivial.
Consequently the character that defines the general length
\(V_{\ell,e}(1)\) is \(1^{-1}\chi_\ell^{-e}=\chi_\ell^{-e}\), and
\(V_{\ell,e}(1)=b_\ell(e)\), including the possibility of infinite
length. This is the equality \texttt{vlen\_\allowbreak{}one}.

Take \(K=K_*(3)\), the imaginary quadratic field of discriminant
\(-3\). Its complex-place count is one and its degree is two. The
Minkowski bound for an integral representative in each ideal class is
\[
   \frac{4}{\pi}\frac{2!}{2^2}\sqrt{3}
   =\frac{2\sqrt{3}}{\pi}<2 .
\]
The representative has positive integral norm and therefore norm one.
It is the unit ideal. This proves \(h(K)=1\).

We verify the local input to the general estimate. Let a valuation on
\(\Omega\) lie over \(3\). Every integer has valuation of absolute value
at most one. If \(\ell\ne3\) and both \(3\) and \(\ell\) had absolute
value less than one, write \(1=3a+\ell b\) by B\'ezout's identity.
The ultrametric inequality would make \(1\) have absolute value less
than one. Hence \(\ell\) is a unit at this valuation. The
\(\ell\)-adic cyclotomic character is trivial on inertia there, as is
the character modulo two. But \(K\) is ramified at \(3\). Thus inertia
at this place cannot act trivially on \(K\) while the whole joint
kernel of the two preceding characters fixes \(K\). This proves the
required noncontainment. The remaining hypotheses of the good-data
predicate follow from \(h(K)=1\), \(\lvert D_K\rvert=3\),
\(\ell\geq5\), \(\ell-1\nmid2e\), and the evenness of \(e\), which gives
\(1=(-1)^e\). The proof uses ramification at \(3\) to test \(K\); \(3\)
is outside the summation range of this polynomial estimate.

Apply the general orbit estimate \(\mathrm{CorGK}\) at this \(K\). Write
\(c_5,a\) for its natural-number constants. Embed \(K\) in \(\Omega\),
take level \(2\), the full subgroup of \((\mathbb Z/2\mathbb Z)^\times\),
and the distinct test primes \(7,11\). Neither divides
\(30\cdot2\lvert D_K\rvert=180\). At each \(\ell\in S\) supply the trivial
character and the good data just verified. Each orbit in the estimate
is the singleton consisting of that character. Its contribution thus
has no orbit multiplicity; the index of the full subgroup is one and
\(\varphi(1)=1\). The finiteness clause of \(\mathrm{CorGK}\) becomes
the claimed finiteness. In its real inequality the sum over characters
is the single term \(b_\ell(e)/r_\ell\). Using
\(\log 7\leq 7\), \(\log 11\leq11\) and expanding the \(a\)-th power of
\(2(e+1)\) gives the claimed bound with
\[
     A=a,\qquad C=2c_5\,2^a(7+11)=36\,c_5\,2^a .
\]
In particular this use of \(7\) is as a test prime, and the assertion
places no requirement that \(7\) belong to \(S\).

The proof of the imported \(\mathrm{CorGK}\) is assembled in
\texttt{Small\allowbreak{}Class\allowbreak{}Numbers.\allowbreak{}Cor\allowbreak{}GKOf}. Starting from \(H'\), the
construction first obtains \(K\)'s size lemma, unramified-coefficient
automorphy from the first endpoint, and cohomological input from the
rank-\(N\) data and the size lemma. It derives a further endpoint from
\texttt{linked\allowbreak{}Unl\allowbreak{}E\_\allowbreak{}alt}. The two algebraic bounds, the auxiliary
algebraic lemma, the Kato input, the two rigidity arguments, the
even-level admissibility result and unramified-coefficient automorphy
give the single-character bound. Combine it with the count of
rank-three unitary representations to obtain \(\mathrm{PropGK}\);
the general theorem \texttt{U3.\allowbreak{}cor\allowbreak{}GK} then yields \(\mathrm{CorGK}\).
Instantiating \(H'\) with the rank-\(N\) patching-and-descent data gives
the unconditional input used in the specialization. This last
paragraph identifies imported results; it does not supply their
individual proofs.
